
\setcounter{page}{2}
\sloppy
\thispagestyle{empty}

\begin{center}
Учреждение образования\\
«Белорусский государственный университет информатики\\
и радиоэлектроники»\\[0.5em]
Факультет информационных технологий и управления
\end{center}

\medskip

\hfill\begin{minipage}{\widthof{Заведующий кафедрой}}


  УТВЕРЖДАЮ

  Заведующий кафедрой

  \tline{\fontsize{10pt}{0pt}{\textit{(подпись)}}}{\hphantom{Заведующий кафедрой}}

  \underline{\hphantom{хорошая дата }}\ \  2026г.

\end{minipage}
\medskip

\begin{center}
  \noindent
  ЗАДАНИЕ

  \noindentпо курсовому проектированию

  \bigskip
  \noindentСтуденту Иосько Михаилу Алексеевичу
\end{center}

\noindent
1. Тема проекта \uline{Автоматизированная информационная система управления интернет-магазина автозапчастей}

\smallskip
\noindent
2. Срок сдачи студентом законченного проекта \uline{\hfill 1 декабря 2026 г. \hfill}

\smallskip
\noindent
3. Исходные данные к проекту \uline{Автоматизация процессов товарного учета автозапчастей, обработки клиентских заказов, контроля движения складских запасов и формирования отчетов по продажам интернет-магазина. Реализация базы данных в СУБД PostgreSQL в третьей нормальной форме без использования объектно-реляционного отображения (ORM) на базе прямых параметризованных SQL-запросов. Реализация каталога номенклатуры с поиском и фильтрацией по категориям и брендам, клиентской корзины, транзакционного оформления заказов со списанием остатков, пакетного обмена номенклатурой в форматах JSON, XML, CSV и генерации аналитической отчетности в форматах PDF, Excel, Word.}

\smallskip
\noindent
4. Содержание расчетно-пояснительной записки (перечень вопросов, которые подлежат разработке):
\begin{description}
  \item \uline{Введение}
  \item \uline{1 Обзор исследуемой области}
  \item \uline{1.1 Описание предметной области}
  \item \uline{1.2 Обзор аналогов}
  \item \uline{1.3 Постановка задачи}
  \item \uline{2 Проектирование системы}
  \item \uline{2.1 Инфологическое и логическое проектирование базы данных}
  \item \uline{2.2 Архитектурное проектирование системы}
  \item \uline{3 Разработка системы}
  \item \uline{3.1 Выбор средств реализации}
  \item \uline{3.2 Реализация базы данных}
  \item \uline{3.3 Программная реализация модулей системы}
  \item \uline{4 Руководство пользователя и администратора}
  \item \uline{Заключение}
  \item \uline{Список использованных источников}
  \item \uline{Приложения}
\end{description}

\smallskip
\noindent
5. Перечень графического материала (с точным обозначением обязательных чертежей и графиков)
\begin{description}
  \item \uline{1. ER-диаграмма в нотации Питера Чена}
  \item \uline{2. Логическая схема базы данных}
  \item \uline{3. Диаграмма вариантов использования (Use Case)}
  \item \uline{4. Диаграмма классов}
  \item \uline{5. Схема алгоритма транзакционного оформления заказа со списанием остатков}
  \item \uline{6. Архитектура развёртывания системы}
  \item \uline{7. Физическая схема базы данных (формат А3)}
\end{description}
\smallskip
\noindent
6. Консультант по проекту (с обозначением разделов проекта) \uline{Трофимович А. Ф. \hfill}

\smallskip
\noindent
7. Дата выдачи задания \uline{\hfill 11 сентября 2026 г.\hfill}

\smallskip
\noindent
8. Календарный график работы над проектом на весь период проектирования (с обозначением сроков выполнения и трудоемкости отдельных этапов):
\begin{description}
  \item \uline{раздел 1 к 24.09 -- 30\%;}
  \item \uline{раздел 2 к 22.10 -- 60\%;}
  \item \uline{раздел 3 к 19.11 -- 90\%;}
  \item \uline{оформление пояснительной записки и графического материала к 11.12 -- 100\%;}
\end{description}

\bigskip
\noindent
\begin{tabular*}{\textwidth}{@{}l@{\extracolsep{\fill}}c@{}r@{}}
  Руководитель\hfill& \tline{(подпись)}{\hphantom{любимый котик}} & \hfill Лаппо А. И.\\
  Задание принял к исполнению\hfill & \tline{(подпись)}{\hphantom{любимый котик}} &\hfill Иосько М. А.
\end{tabular*}
